\section{Conclusion}
\label{sec:conclusion}

We have presented a controlled study of prompt sensitivity across six tasks, a
held-out seventh, eight models (1B--frontier), and 21,858 evaluations, scored
on deterministic ground truth wherever it exists.
Through additive prompt levels, a randomised layer-ordering ablation, and an
irrelevant-padding control, we show that response quality is governed by
\textbf{information sufficiency}, and that for structured tasks the
task-critical information is the \emph{output schema}.

Our central finding is \textbf{schema gating}: a single schema-resolving layer
captures 87--95\% of all quality gain for classification, product extraction,
and a held-out NER task, with everything else (definitions, personas, worked
examples) adding little.
This concentration is sharpest on objective ground truth; the LLM judge had
masked it, illustrating why metric choice matters.
The remaining tasks fall into two further classes distinguished by mechanism:
\textbf{difficulty-gated} (instruction following, whose sensitivity grows
$6.6\times$ with harder examples) and \textbf{knowledge-gated} (QA, math,
summarisation, at ceiling from the bare input).
Because we predicted the held-out task's class and gain locus before running,
the taxonomy has out-of-sample support, not merely descriptive fit.
The practical message is simple: for structured tasks, invest the prompt budget
in one well-formed schema specification and stop; for knowledge tasks, the bare
question suffices.

\paragraph{Future directions.}
The schema-gating mechanism makes a sharp prediction we can test out of sample:
any task whose output the model cannot infer from the input (code generation
with a type signature, data-to-text with a template, dialogue-state tracking
with a slot schema) should concentrate its gain at the format layer.
A held-out battery of such tasks would validate the difficulty- and
knowledge-gated classes as the NER task validated the schema-gated one.
Mechanistic-interpretability analyses could directly test how a model
internalises the format layer. Agentic settings, where a tool-call schema is
exactly the kind of format specification our mechanism predicts should
dominate, would extend the finding to long-horizon deployments.
